% UNINTEGRATED WORKING APPENDIX FRAGMENT, 24 October 2026.
% Same general compact paper family; no standalone manuscript or compilation.
% Contains the reference and local isolated-gain bodies; full finite-pair
% and whole-A18 compositions are still to be transcribed.

\section{The whole-\(A_{18}\) endpoint argument}
\label{app:wholeA18}
Here the entire diagram is \(A_{18}\), with no spectators.
We keep one original farthest vector in the zero-sum
subspace of \(\mathbb R^{19}\).
Write \(D_0=\omega_1^2=\omega_{18}^2=36/19\).
The endpoint opposites are actual distinct walls.

\subsection{The reference under the full norm bound}
The \(N=19\) central inventory and helper nonadjacency
above precede any contrary energy assumption.
Its near-four roles can also be excluded under
only \(\norm x<1\), using their sharper actual mates.
For a left near-four at \(i=5,6\), release \(j=i-1\).
The positive replacement has diagram
\(B_{20-i}+A_{i-2}\) and cosine
\[
 C_j=\sqrt{2(i-1)/19}
     \bigl((20-i)/2+(19-i)/\sqrt2\bigr)>11.
\]
Its normalized height exceeds \(\sqrt{5/6}>.91\).
At the other end a mate or a strong \(13/14\)
has demand \(>209/270>.774\).
Their correlation is at most
\(\sqrt{15/91}<.41\).
Cauchy contradicts \(1.684^2=2.835856>2.82\).
Thus all central weak sites are nonadjacent pure
helpers already under the full norm bound.

Use the same \(N=19\) rational heights and chord
omissions. If \(k\) is the number of weak end
sites \(5,14\), one reference \(w\) has energy
\[
 E_0\geq.9240352,\qquad E_1>.8937,\qquad E_2>.86414.
\]
Its left and right block lengths \(m\) are \(5\)
or \(6\), with signed means
\[
 b_5=H_5/(5\sqrt2)\in(.29,.298),\qquad
 b_6=H_6/(6\sqrt2)\in(.25,.266).
\]
The separate boundary dual prices satisfy
\[
 a_5\geq(.4208-.151)/2=.1349,\qquad
 a_6\geq(451/1200-.108)/2>.1339.
\]
Every omission is included and can only increase
the adjacent available price.

For an actual endpoint pressure \(\ip{x}{V}\geq K\),
subtract \(B\), a positive combination of boundary
duals, so that \(Y=V-B\) has zero mean on each
end block. If the individual debit \(\mu B\)
fits, one full square gives gain
\[
 2\mu(K-\ip{w}{B})-\mu^2Y^2.
\]
Opposite-block residuals have zero cross.
For a coordinate pressure \(x_{\rm end}\geq t\),
the legally optimized reference gain is
\(m(t-b_m)^2/(m-1)\).
For all thresholds used below, the debit is
less than \(.061\), below each separate price.
If either coordinate were \(\geq.603\), the
three energy bounds plus this gain would exceed
\(1\); the tight case is
\[
 .86414+(6/5)(.603-.266)^2=1.0004228.
\]
Therefore \(x_1,-x_{19}<.603\).
The genuine endpoint-edge inequality and
\(.603^2<90/247=(18/19)(5/13)\) imply
\begin{equation}\label{app:eq:A18shortends}
 C_L,C_R<9/4.
\end{equation}
This bounds actual edge cosines, not cap owns.

\subsection{Complete local isolated-endpoint gains}
An isolated left endpoint has
\(c=A p-g\omega_1,\ A^2=D_0g^2-1>0\).
Own zero, \(3,4\) has nonpositive time; own
\(6\) and parallel have the positive actual
nulls \(\sqrt3c+2n_1+n_2\) and \(c+n_1\).
The full remaining list is finite \(5\), finite
\(7\), finite at least \(8\), or dotted.
For a pressure normalized to raw coordinate
sum \(\sqrt2\), subtract \(\omega_m/m\);
the residual is zero-sum on the relevant block.
Every price below is paid before applying its
actual pressure. The resulting complete gain
table is
\[
\begin{array}{c|c|c|c}
\text{own}&\mu&\text{gain on }m=5>&
 \text{gain on }m=6>\\ \hline
5&.6&.02904&.059\\
7&.4&.05152&.07677\\
\text{finite }\geq8&1/4&.0471&157/2400\\
\text{dotted}&1/4&.06&4699/60000.
\end{array}
\]
We prove every row, including continuous dotted
owns and the unbounded finite-label tail.

For dotted \(g>1\), put \(u=g^{-2}\in(0,1)\).
The actual positive core \(c+n_1/g\) gives
\[
 V=\omega_1-un_1,\quad
 K(u)=\sqrt{D_0-u}-\sqrt{u(1-u)},\quad
 Y^2=2-2/m-2u+u^2.
\]
It has square \(u-1<0\), positive time,
zero cap product and negative old-\(1\)
product; positivity checks every unknown wall.
At \(\mu=1/4\) the gain is
\[
 G_m(u)=\tfrac12(K(u)-H_m/m)
        -(2-2/m-2u+u^2)/16.
\]
Its shift is exactly
\(G_6-G_5=1099/60000\).
To prove \(G_5>.06\) uniformly, it suffices that
\[
 K(u)>.7408-u/4+u^2/8.
\]
The upper bound
\(\sqrt{u(1-u)}\leq1/2-(u-1/2)^2\)
follows by positive squaring, with residual
\((u-1/2)^4\).
Use lower radical values \(1.376,1.1809,1.0699,.9459\)
at \(u=0,1/2,3/4,1\); their squares are
\(1.893376,1.39452481,1.14468601,.89472681\),
strictly below \(D_0-u\).
Concavity places \(\sqrt{D_0-u}\) above the
successive chords. The three resulting
strict quadratic lower bounds are
\[
\begin{array}{c|l}
[0,1/2]&.3852-1.1402u+(7/8)u^2\\
[1/2,3/4]&
(7/8)(u-597/875)^2+8357/1750000\\
[3/4,1]&.4511-1.246u+(7/8)u^2.
\end{array}
\]
The first decreases to \(677/20000>0\);
the last increases from \(703/80000>0\).
This proves the entire dotted row.

For finite \(g<1\), put \(u=g^{-2}\), \(d=D_0-1=17/19\).
The actual finite projection onto \(c,n_1\)
gives
\[
 V=\omega_1-n_1,\quad
 K_f=\frac{d-\sqrt{d(u-1)}}{\sqrt{D_0-u}},\quad
 Y^2=1-2/m.
\]
Writing \(s=\sqrt{u-1}\),
\(K_f^2=d(\sqrt d-s)/(\sqrt d+s)\), increasing
with \(g\).
For every native own at least \(8\),
\(s\leq\sqrt2-1<.414214\).
Using \(.8947<d<.8948,\sqrt d>.9459\),
\[
 (d-.59^2)\sqrt d>.51702894,\qquad
 (d+.59^2)s<.5148265806.
\]
Hence \(K_f>.59\). Quarter-payment gives
\(G_{f,5}>.0471\) and
\(G_{f,6}>.0471+1099/60000=157/2400\).

For native \(5\), let \(u=6-2\sqrt5\) and
\(b=D_0-3/2=15/38\).
Its actual \(H_4\) correction
\(\eta=c+g(3n_1/2+n_2+n_3/2)\)
has square \(\delta=1-3g^2/2>0\).
Use the actual finite point
\(p+(A/\delta)\eta\), not \(\eta\).
The pressure is
\[
 V=\omega_4/4,\qquad
 K_5=\frac{b-\sqrt{b(u-3/2)}}{\sqrt{D_0-u}},
 \qquad Y^2=1/2-2/m.
\]
For \(s=\sqrt{u-3/2}<.167\),
\(.3947<b<.3948,\sqrt b>.6282\) give
\[
 (b-.475^2)\sqrt b>.106212915,\qquad
 (b+.475^2)s<.103610975.
\]
Thus \(K_5>.475\).
The payment \(.6/m\leq.12\) fits each
boundary coefficient and gives
\[
 G_{5,5}>1.2(.475-.4208)-.36(.1)=.02904,\qquad
 G_{5,6}>1.2(.475-451/1200)-.36/6=.059.
\]

For native \(7\),
\(.900968<g<.900969\) and \(1.231<u=g^{-2}<1.232\).
The positive actual core
\(c+g^{-1}n_1+2(g^{-1}-g)n_2\), divided by \(g\),
gives
\[
 V=\omega_1-un_1-2(u-1)n_2,\quad
 K_7=\sqrt{D_0-u}-\sqrt{7u-3u^2-4},
\]
\[
 Y^2=6-8u+3u^2-2/m.
\]
Its selected products are zero on \(c,n_1\),
negative on \(n_2,n_3\), and nonpositive on
every unknown wall; time is positive.
On the stated interval,
\(D_0-u>.814^2\) and
\(0<7u-3u^2-4<.070917<.2664^2\), so \(K_7>.545\).
The residual bounds are \(Y_5^2<.299\),
\(Y_6^2<.366\).
Payment \(.4/m\leq.08\) gives gains
\[
 .8(.545-.4208)-.16(.299)=.05152,\qquad
 .8(.545-451/1200)-.16(.366)>.07677.
\]
This proves the table and all its actual-wall
and individual-price hypotheses.
\subsection{The complete two-finite endpoint exclusion}
\label{app:A18finitepair}
Suppose both isolated endpoint owns are finite.
For \(D=36/19,e=2/19\), their actual mutual
magnitude is
\[
 M(g,h)=\sqrt{Dg^2-1}\sqrt{Dh^2-1}-egh.
\]
It increases strictly in each own for
\(g,h\geq g_5=\cos(\pi/5)\).
For example the derivative sign follows by
squaring positive terms from
\(g/\sqrt{Dg^2-1}>1/\sqrt D\) and
\((D^2-e^2)h^2>D\), since
\(D^2-e^2=68/19\) and \(g_5^2>9/17\).

If either own is \(5\), monotonicity gives
\[
 0<\frac{17\sqrt5-25}{76}=M(g_5,g_5)
 \leq M(g_5,h)<M(g_5,1)
 <\sqrt{(9\sqrt5-11)/38}<1/2.
\]
That is forbidden for distinct actual walls.
Thus both own labels \(m,n\) are at least \(7\).
Relabel so \(m\leq n\). Now \(M>0\) already
at \(m=n=7\), while
\(M<M(1,1)=15/19<\cos(\pi/5)\).
The full mutual alphabet therefore leaves
exactly native \(3\) or native \(4\).
Put
\[
 H(g,h)=272g^2h^2-144(g^2+h^2)-8gh+57,
\]
\[
 F_S(g,h)=136g^2h^2-72(g^2+h^2)+19-4Sgh.
\]
They have the signs of \(M-1/2\) and
\(M-1/\sqrt2\), respectively, with
\(S=\sqrt2\). This follows by the
difference of squares against \(egh\)
plus the target; the other factor is
positive. Both increase in each argument
on \([.9,1]^2\), for \(1.4\leq S\leq1.5\),
and \(F_S\) decreases with \(S\).

For mutual \(3\), \(m\geq8\) gives
\[
 M(g_8,g_8)=(17\sqrt2-4)/38>1/2.
\]
At \(m=7\) the exact rational native
enclosures below give
\[
 M(g_7,g_8)<1/2<M(g_7,g_9).
\]
At the upper \(7/8\) pair the bounds
\(u<.8117452,v<.8535543,\ u+v>1.6652993\),
\(uv<.692869,gh>.832387\) yield
\(H<-10018272/10^7<0\).
At the lower \(7/9\) pair,
\(u>.8117433,v>.8830210,\ u+v<1.6947645\),
\(uv>.716786,gh<.846633\) yield
\(H>1146640/10^6>0\).
Thus no integer \(n\geq m\) gives mutual \(3\).

For mutual \(4\), if \(m\leq10\), increase
\(g\) to the upper \(g_{10}^+=.951057\)
and \(h\) to \(1\). Since
\((g_{10}^+)^2<.9045095\) and
\(4(1.414213)g_{10}^+>5.37997\),
\[
 F_{\sqrt2}(g,h)
 \leq F_{1.414213}(g_{10}^+,1)<-.491362<0.
\]
If \(m\geq15\), monotonicity gives
\[
 M\geq(34/19)\cos^2(\pi/15)-1>1/\sqrt2.
\]
Use \(\cos(\pi/15)>10783/11025\),
from \(\cos x>1-x^2/2,\pi<22/7\), and
\(\sqrt2<99/70\). The rational remaining
margin is
\[
 (34/19)(10783/11025)^2-1-99/140
   =1498515515/323324662500>0.
\]
Only \(m=11,12,13,14\) remain.
Their complete consecutive-integer gaps are
\[
\begin{array}{c|c|c}
m&\text{last }n\text{ below target}&
 \text{first }n\text{ above target}\\ \hline
11&31&32\\12&20&21\\13&16&17\\14&15&16.
\end{array}
\]

Here is full rational arithmetic for these gaps.
For \(g=g_m,S=\sqrt2\), set
\[
 A=136g^2-72,\quad B=4Sg,\quad
 C=64g^2-53-B,\quad T=2A-B.
\]
At \(h=1-\delta\),
\(F_S(g,h)=C-T\delta+A\delta^2\).
The native enclosures give
\[
\begin{array}{c|cc|cc}
m&g^2\text{ lower}&g^2\text{ upper}&
 B\text{ lower}&B\text{ upper}\\ \hline
11&.9206248&.9206269&5.42770&5.42773\\
12&.9330110&.9330131&5.46409&5.46412\\
13&.9427264&.9427284&5.49246&5.49249\\
14&.9504826&.9504847&5.51501&5.51504.
\end{array}
\]
Thus the following looser bounds hold:
\[
\begin{array}{c|cc|cc|c}
m&C_-&C_+&T_-&T_+&A_+\\ \hline
11&.49225&.49243&100.981&100.984&53.2053\\
12&1.24858&1.24875&104.314&104.317&54.8898\\
13&1.84199&1.84216&106.928&106.931&56.2111\\
14&2.31584&2.31602&109.016&109.018&57.2660.
\end{array}
\]
For last-below \(n\), use its upper cosine
and \(S=1.414213\).
The \(\delta\) numerators over \(10^6\) are
\(5130,12311,19214,21852\);
upper numerators of \(A\delta^2\) are
\(1401,8400,20800,27374\), also over \(10^6\).
The exact resulting upper numerators over
\(10^9\) are
\[
 -24201530,\ -27059654,\ -191554592,\ -38823632.
\]
For first-above \(n\), use its lower cosine
and \(S=1.414214\).
The \(\delta\) numerators are
\(4816,11170,17027,19215\);
dropping the positive quadratic term gives
lower numerators over \(10^9\)
\[
 5911056,\ 83359110,\ 21275863,\ 221059130.
\]
All signs are strict.
Monotonicity in \(g,h\) proves every integer
\(n\geq m\) lies on one side of the stated
gap, so mutual \(4\) is impossible as well.
This proves the two-finite exclusion.

\subsection{The finite rational angle certificate}
\label{app:A18anglecertificate}
For completeness, the only native enclosures
used in the preceding finite calculation are:
\[
\begin{array}{c|r|r@{\qquad}c|r|r}
j&10^6g_j^-&10^6g_j^+&
j&10^6g_j^-&10^6g_j^+\\ \hline
7&900968&900969&8&923879&923880\\
9&939692&939693&10&951056&951057\\
11&959492&959493&12&965925&965926\\
13&970941&970942&14&974927&974928\\
15&978147&978148&16&980785&980786\\
17&982973&982974&20&987688&987689\\
21&988830&988831&31&994869&994870\\
32&995184&995185&&&
\end{array}
\]
These are exact rational inequalities, not
floating-point evaluations.
Machin's identity
\(\pi=16\arctan(1/5)-4\arctan(1/239)\)
follows by the tangent addition formula
and its first-quadrant sign.
For
\[
 A_s=\sum_{r=0}^s\frac{(-1)^r}{(2r+1)5^{2r+1}},
 \qquad
 B_s=\sum_{r=0}^s\frac{(-1)^r}{(2r+1)239^{2r+1}},
\]
alternating remainders and rational
cross multiplication give
\[
 31415926/10^7<16A_5-4B_2<\pi
 <16A_6-4B_1<31415927/10^7.
\]
For \(0<x<.45\), the alternating cosine
Taylor sums give \(C_{10}(x)<\cos x<C_8(x)\).
Put \(a_j=10^7j,p_-=31415926,p_+=31415927\), and
\[
\begin{split}
Q_{10}(a,p)={}&3628800a^{10}-1814400p^2a^8
 +151200p^4a^6-5040p^6a^4\\
 &+90p^8a^2-p^{10},\qquad Q_8=Q_{10}+p^{10}.
\end{split}
\]
For every displayed row with integers \(L_j,U_j\),
the finite exact certificate is
\[
 10^6Q_{10}(a_j,p_+)-3628800L_ja_j^{10}>0,
\]
\[
 3628800U_ja_j^{10}-10^6Q_8(a_j,p_-)>0.
\]
There are exactly thirty integer inequalities;
their positivity is obtained by direct
integer multiplication of the specified
finite expressions.
The alternating remainder signs and cosine
monotonicity then prove all fifteen rows.
This is a terminating exact verification
recipe with a complete stated output;
no execution of a checker is claimed by
this manuscript transcription.
The historical unused cyclotomic fallback
is not an input.
Finally \(1.414213<\sqrt2<1.414214\)
follows by squaring those two rationals.

\subsection{The complete both-isolated endpoint conclusion}
The local isolated gains, with their separate
prices and disjoint residual supports, exclude
two dotted endpoints:
the three reference-plus-gain lower bounds are
\[
 .9240352+2(.06)>1,\quad
 .8937+.06+4699/60000>1,\quad
 .86414+2(4699/60000)>1.
\]
For one finite and one dotted endpoint, the
full four-pattern totals are
\[
\begin{array}{c|c|c|c|c}
\text{finite own}&5/5&6/5&5/6&6/6\\ \hline
\geq8&1.0311352&>1.01911666&>1.01911666&>1.00787333\\
5&1.0130752&1.0127&>1.00105666&>1.00145666\\
7&1.0355552&>1.03047&>1.02353666&>1.01922666.
\end{array}
\]
In the middle two columns the block order
is finite/dotted. The gains add within
one square because their residuals occupy
opposite zero-sum blocks.
Together with the two-finite theorem,
these exhaust every isolated endpoint own.
Thus both actual endpoint opposites cannot
be isolated.




